\documentclass[a4paper]{article}
\usepackage[affil-it]{authblk}
\usepackage[english]{babel}
\usepackage[backend=bibtex,style=numeric]{biblatex}

\usepackage{geometry}
\geometry{margin=1.5cm, vmargin={0pt,1cm}}
\setlength{\topmargin}{-1cm}
\setlength{\paperheight}{29.7cm}
\setlength{\textheight}{25.3cm}

% useful packages.
\usepackage{amsfonts}
\usepackage{amsmath}
\usepackage{amssymb}
\usepackage{amsthm}     
\usepackage{enumerate} %enumerate environment
\usepackage{graphicx}
\usepackage{multicol}
\usepackage{fancyhdr}
\usepackage{layout}
\usepackage{ctex}      %中文环境
\usepackage{booktabs}  %三线表
\usepackage{verbatim}  %comment environment
\usepackage{float}     %强制表格图片位置
\usepackage{listings}  %insert code in latex
\usepackage{fontspec}  %font support
\usepackage{tikz}
\usepackage{makecell}
\usetikzlibrary{intersections}
\usepackage{arydshln}

\usepackage[colorlinks,linkcolor=blue]{hyperref}
\usepackage{xcolor}
\definecolor{mygreen}{rgb}{0,0.6,0}
\definecolor{lightgrey}{rgb}{0.95,0.95,0.95}
\definecolor{winered}{rgb}{0.5,0,0}
\definecolor{mymauve}{rgb}{0.58,0,0.82}
\lstset{
      backgroundcolor=\color{lightgrey},
      basicstyle             = \ttfamily,
      breakatwhitespace      = false,
      breaklines             = true,
      captionpos             = none,
      commentstyle           = \color{mygreen}\bfseries,
      extendedchars          = false,
      keepspaces             = true,
      keywordstyle           = \color{winered},         % keyword style
      language               = C++,                     % the language of code
      otherkeywords          = {string},
      numberstyle            = \tiny\color{mygray},
      rulecolor              = \color{black},
      showspaces             = false,
      showstringspaces       = false,
      showtabs               = false,
      stepnumber             = 1,
      stringstyle            = \color{mymauve},         % string literal style
      tabsize                = 2,
      title                  = \lstname,
      aboveskip              = 0pt,
      belowskip              = 0pt,
      columns                = flexible
}

% some common command
\newcommand{\dif}{\mathrm{d}}
\newcommand{\avg}[1]{\left\langle #1 \right\rangle}
\newcommand{\norm}[1]{\left\| #1 \right\|}
\newcommand{\difFrac}[2]{\frac{\dif #1}{\dif #2}}
\newcommand{\pdfFrac}[2]{\frac{\partial #1}{\partial #2}}
\newcommand{\OFL}{\mathrm{OFL}}
\newcommand{\UFL}{\mathrm{UFL}}
\newcommand{\fl}{\mathrm{fl}}
\newcommand{\op}{\odot}
\newcommand{\Eabs}{E_{\mathrm{abs}}}
\newcommand{\Erel}{E_{\mathrm{rel}}}
\addbibresource{citation.bib}

\begin{document}
% =================================================
\title{Numerical Analysis chapter 2 theoretical homework}

\author{Zhouyue Qian 3220104074
  \thanks{Electronic address: \texttt{3220104074@edu.zju.cn}}}
\affil{(Information and Computational Science Class 2201), Zhejiang University }


\date{Due time: \today}

\maketitle

%\begin{abstract}
%    The abstract is not necessary for the theoretical homework, 
%    but for the programming project, 
%    you are encouraged to write one.      
%\end{abstract}





% ============================================
\section*{I}
\subsection*{Question1}
\textit{Solution:}
The linear interpolation of $f(x)$ at $x_0$ is given by $$ p_1(f; x) = f(x_0) + \frac{f(x_1) - f(x_0)}{x_1 - x_0} (x - x_0). $$
As $f(x) = \frac{1}{x}$,$x_0 = 1$,$x_1 = 2$ we have:
$$f(x_0) = f(1) = 1,$$
$$f(x_1) = f(2) = \frac{1}{2}. $$
So the linear interpolation of $f(x)$ at $x_0 = 1$ is:
$$p_1(f; x) = 1 + \frac{\frac{1}{2} - 1}{2 - 1} (x - 1) = 1 - \frac{1}{2}(x - 1) = \frac{3 - x}{2}.$$
The error is   \[f(x) - p_1(f; x) = \frac{f''(\xi(x))}{2} (x - 1)(x - 2).\]
As $f(x) = \frac{1}{x}$, $f''(x) = 2x^{-3}$, \[ \frac{1}{x} - \frac{3 - x}{2} = \frac{\frac{2}{\xi(x)^3}}{2} (x - 1)(x - 2). \]
\[ \xi(x)=\sqrt[3]{2x } . \]\hfill $\Box$ \cite*{deepseek}

\subsection*{Question2}
\textit{Solution:}
$$\max \xi(x) = \xi(2) = \sqrt[3]{4} , \quad \min \xi(x) = \xi(1) = \sqrt[3]{2}, \quad \max f''( \xi(x) ) = 1$$

\section*{II}
First, we use Lagrange interpolation to construct a polynomial \( L(x) \) such that \( L(x_i) = f_i \) for \( i = 0, 1, \ldots, n \). The Lagrange interpolation polynomial can be expressed as:
\[
L(x) = \sum_{i=0}^n f_i \ell_i(x)
\]
where \(\ell_i(x)\) are the Lagrange basis functions:
\[
\ell_i(x) = \prod_{\substack{0 \leq j \leq n \\ j \neq i}} \frac{x - x_j}{x_i - x_j}
\]


If $p(x) = L(x)$, adn $\forall x , p(x)>0$, then $p(x)$ is what we want. If not, we could add another polynomial $q(x) = A\cdot \prod_{i=1}^{n}(x-x_i)^{2},A>0$.
By adjusting the constant \( A \), it is possible to ensure that \( q(x) + L(x) > 0 \) for all \( x \), because \( q(x) \) is always greater than 0.
And this function does not exceed the 2n-th order.\hfill $\Box$

\section*{III}
\subsection*{Question1}

\textbf{Base Case (n=1):}

\[
f[t, t+1] = \frac{f(t+1) - f(t)}{(t+1) - t} = \frac{e^{t+1} - e^t}{1} = e^t (e - 1).
\]

This matches \( \frac{(e-1)^1}{1!} e^t \).

\textbf{Inductive Step:}
Assume the statement holds for some \( k \geq 1 \), i.e.,
\[
f[t, t+1, \ldots, t+k] = \frac{(e-1)^k}{k!} e^t.
\]

We need to prove it for \( k+1 \):
\[
f[t, t+1, \ldots, t+k, t+k+1] = \frac{(e-1)^{k+1}}{(k+1)!} e^t.
\]

\[
f[t, t+1, \ldots, t+k, t+k+1] = \frac{f[t+1, t+2, \ldots, t+k+1] - f[t, t+1, \ldots, t+k]}{(t+k+1) - t}.
\]

\[
f[t+1, t+2, \ldots, t+k+1] = \frac{(e-1)^k}{k!} e^{t+1},
\]
\[
f[t, t+1, \ldots, t+k] = \frac{(e-1)^k}{k!} e^t.
\]

Thus:
\[
f[t, t+1, \ldots, t+k, t+k+1] = \frac{\frac{(e-1)^k}{k!} e^{t+1} - \frac{(e-1)^k}{k!} e^t}{1} = \frac{(e-1)^k}{k!} (e^{t+1} - e^t) = \frac{(e-1)^k}{k!} e^t (e - 1).
\]

Simplifying:
\[
f[t, t+1, \ldots, t+k, t+k+1] = \frac{(e-1)^{k+1}}{(k+1)!} e^t.
\]\hfill $\Box$
\subsection*{Question2}
\[
f[0, 1, \ldots, n] = \frac{(e-1)^n}{n!} e^0 = \frac{(e-1)^n}{n!}.
\]

Thus:
\[
\frac{(e-1)^n}{n!} = \frac{1}{n!} f^{(n)}(\xi).
\]

Since \( f^{(n)}(x) = e^x \):
\[
(e-1)^n = e^\xi.
\]

Taking the natural logarithm:
\[
\xi = \ln((e-1)^n) = n \ln(e-1).
\]

So, \(\xi\) is located to the right of \( n/2 \).
\section*{IV}
\subsection*{Question1}
\[
\begin{array}{c|cccc}
x & f(x) & f[x, x+1] & f[x, x+1, x+2] & f[x, x+1, x+2, x+3] \\
\hline
0 & 5 & & & \\
1 & 3 & \frac{3 - 5}{1 - 0} = -2 & & \\
3 & 5 & \frac{5 - 3}{3 - 1} = 1 & \frac{1 - (-2)}{3 - 0} = 1 & \\
4 & 12 & \frac{12 - 5}{4 - 3} = 7 & \frac{7 - 1}{4 - 1} = 2 & \frac{2 - 1}{4 - 0} = \frac{1}{4} \\
\end{array}
\]\cite*{deepseek}

\[
p_3(f; x) = f(0) + f[0, 1](x - 0) + f[0, 1, 3](x - 0)(x - 1) + f[0, 1, 3, 4](x - 0)(x - 1)(x - 3)
\]

Simplifying:
\[
p_3(f; x) = 5 - \frac{9}{4}x  + \frac{1}{4} x^3.
\]
\subsection*{Question2}
\[
p_{3}^{'}(f; x) = - \frac{9}{4} + \frac{3}{4} x^2.
\]
so, $x = \sqrt{3} $, when the derivative of the polynomial is zero. \hfill $\Box$
\section*{V}
\subsection*{Question1}
We assume \[ p_5(f; x) = \sum_{i=0}^{5} a_i x^i \], and we have:
\[ p_5(0) = f(0), \quad p_5(1) = f(1), \quad p_5'(1) = f'(1), \quad p_5''(1) = f''(1), \quad p_5(2) = f(2), \quad p_5'(2) = f'(2) \]
So, $a_0 = 0$
\begin{equation}
  \begin{pmatrix}
  1 & 1 & 1 & 1 & 1 \\
  1 & 2 & 3 & 4 & 5 \\
  0 & 2 & 6 & 12 & 20 \\
  2 & 4 & 8 & 16 & 32 \\
  1 & 4 & 12 & 32 & 80
  \end{pmatrix}
  \begin{pmatrix}
  a_1 \\
  a_2 \\
  a_3 \\
  a_4 \\
  a_5
  \end{pmatrix}
  =
  \begin{pmatrix}
  1 \\
  7 \\
  42 \\
  128 \\
  448
  \end{pmatrix}
\end{equation}
So, we have: $f[0,1,1,1,2,2] = a_5 = 30$

\subsection*{Question2}
$f[0,1,1,1,2,2] = a_5 = 30 = \frac{f^{(5)(\xi )}}{5!}$
So, $\xi = \sqrt{\frac{10}{7}} $ \hfill $\Box$

\section*{VI}
\subsection*{Question1}
\[
  \begin{array}{c|cccc}
  0 & 1 & & & \\
  1 & 2 & 1 & & \\
  1 & 2 & -1 & -2& \\
  3 & 0 & -1 & 0 & \frac{2}{3} \\
  3 & 0 & 0 & \frac{1}{2} & \frac{-5}{36} \\
  \end{array}
  \]
  $$p_4(f;x) = 1 + 1*(x-0)+ -2*(x-0)(x-1) + \frac{2}{3}*(x - 0)(x - 1)^{2} + \frac{-5}{36}(x - 0)(x - 1)^{2}(x - 3)$$ 

$p_4(f;2)=1+2-2*2 + \frac{2}{3}*2-\frac{5}{36}*2*(-1)=\frac{11}{18}$

\subsection*{Question2}
\[
  \exists \xi \in (0, 3), \text{ such that } \left| f(x) - p_4(f; x) \right| = \left| \frac{f^{(5)}(\xi)}{5!} (x - 0)(x - 1)^2 (x - 3)^2 \right| \leq \frac{M}{120} \left| (x - 0)(x - 1)^2 (x - 3)^2 \right|
\]

$x = 2$ , $R_4(2)\leqslant \frac{M}{60} $

\section*{VII}
\subsection*{forwards}
\[
\Delta f(x) = f(x + h) - f(x),
\]
\[
\Delta^{k+1} f(x) = \Delta \Delta^k f(x) = \Delta^k f(x + h) - \Delta^k f(x).
\]

We use mathematical induction to prove the formula for the forward difference.

\textbf{Base case:} When \( k = 1 \),
\[
\Delta f(x) = f(x + h) - f(x).
\]

\textbf{Inductive hypothesis:} Assume the formula holds for some \( k \), i.e.,
\[
\Delta^k f(x) = k! h^k f[x_0, x_1, \ldots, x_k].
\]
\[
\Delta^k f(x + h) = k! h^k f[x_1, x_2, \ldots, x_{k+1}].
\]

\textbf{Inductive step:} Prove the formula holds for \( k+1 \)

\[
\Delta^{k+1} f(x) = \Delta (\Delta^k f(x)).
\]

\[
\Delta^{k+1} f(x) = \Delta^k f(x + h) - \Delta^k f(x) = k! h^k f[x_1, x_2, \ldots, x_{k+1}] - k! h^k f[x_0, x_1, \ldots, x_k].
\]
\[
f[x_0, x_1, \ldots, x_{k+1}] = \frac{f[x_1, x_2, \ldots, x_{k+1}] - f[x_0, x_1, \ldots, x_k]}{x_{k+1} - x_0}.
\]
\[
\Delta^{k+1} f(x) = k! h^k \cdot \frac{f[x_1, x_2, \ldots, x_{k+1}] - f[x_0, x_1, \ldots, x_k]}{h} = (k+1)! h^{k+1} f[x_0, x_1, \ldots, x_{k+1}].
\]
\[
  \Delta^{k+1} f(x) = (k+1)! h^k \cdot \frac{f[x_1, x_2, \ldots, x_{k+1}] - f[x_0, x_1, \ldots, x_k]}{(k+1)h} = (k+1)! h^{k+1} f[x_0, x_1, \ldots, x_{k+1}].
\] \hfill $\Box$

\subsection*{backwards}
Similarly, we can prove the formula for the backward difference.

\section*{VIII}
\begin{align*}
  f[x_0, x_0, x_1, \ldots, x_n] &= \lim_{x \to x_0} f[x_0, x_0, x_1, \ldots, x_n, x] \\
                                &= \lim_{x \to x_0} \frac{f[x_1, \ldots, x_n, x] - f[x_0, x_1, \ldots, x_n]}{x - x_0} \\
                                &= \frac{\partial}{\partial x_i} f[x_0, x_1, \ldots, x_n]
\end{align*}
Similarly, we can prove the formula for the general case.
\[
\frac{\partial}{\partial x_i} f[x_0, x_1, \ldots, x_n] = f[x_0, \ldots, x_i, x_i, \ldots, x_n]
\]

\section*{IX}
Let $g(x)=\frac{f(x)}{a_0}$, so we have 
\[
  \max _{(a,b)}g(x)\ge \frac{1}{2^{n-1}} (\frac{b-a}{2})^{n}.
\]
So, $\min  \max f(x) = a_0\frac{1}{2^{n-1}} (\frac{b-a}{2})^{n}.$

\section*{X}
Since $T_n(x)$ is Chebyshev polynomial, assuming $T_n(x) = cosh(n*arccosh(x))$, it has $n+1$ extremum points.
As we all know, $\|T_n(x)\|_{\infty} = 1$.
Let $Q(x) = \frac{T_n(x)}{T_n(a)} - p(x)$, we have $$ Q(a) = \frac{T_n(a)}{T_n(a)} - p(a) = 0.$$
Now, we assume that the proposition to be proved is not true.
$$\exists p \in \mathbb{P},\text{ such that } \max_{x \in [-1,1]} |p(x)| < \frac{1}{T_n(a)}  .$$
$$Q(x_k^{'}) = \frac{T_n(x_k^{'})}{T_n(a)} - p(x_k^{'}), k = 0,1, \ldots n .$$
\text{Assume that } Q(x) \text{ has } n+1 \text{ extremum points where the signs of the two neighboring values are different, then } Q(x) \text{ has at least } n \text{ zeros.}
Moreover, we have $x = 1$ as a zero point. So, we have at least $n+1$ zeros. However, \text{The degree of } Q(x) \text{ is } n, which means it has n zeros most.
Therefore, $Q(x) \equiv  0$,$\max_{[-1,1]} p(x) = \frac{1}{T_n(x)}$, which contradicts the assumption.  \hfill $\Box$

\section*{XI}
\textbf{Proof:}
\[
b_{n,k}(t) = \binom{n}{k} t^k (1-t)^{n-k}.
\]
We need to prove:\[
b_{n-1,k}(t) = \frac{n-k}{n} b_{n,k}(t) + \frac{k+1}{n} b_{n,k+1}(t).
\]
First, we have:
\[
b_{n-1,k}(t) = \binom{n-1}{k} t^k (1-t)^{n-1-k}.
\]
\[
b_{n,k}(t) = \binom{n}{k} t^k (1-t)^{n-k}.
\]
\[
b_{n,k+1}(t) = \binom{n}{k+1} t^{k+1} (1-t)^{n-k-1}.
\]

Now we need to prove:
\[
  \binom{n-1}{k} = \frac{n-k}{n}\binom{n}{k}(1-t)+\frac{k+1}{n}\binom{n}{k+1}t
\]

\begin{align*}
  \binom{n-1}{k} &= \frac{n-k}{n}\binom{n}{k}(1-t)+\frac{k+1}{n}\binom{n}{k+1}t\\
  &=\frac{n-k}{n}\frac{n!}{k!(n-k)!}(1-t)+\frac{k+1}{n}\frac{n!}{(k+1)!(n-k-1)!}t\\
  &=\frac{(n-k)(k+1)n!+(k+1)n!(n-k)}{n(k+1)!(n-k)!}\\
  &=\frac{(n-1)!}{k!(n-k-1)!}
\end{align*}
The proof is completed.

\section*{XII}
\textbf{Proof:}
\begin{align*}
  \int_{0}^{1} \binom{n}{k}t^{k}(1-t)^{n-k}dt &= \binom{n}{k}\int_{0}^{1}t^{k}(1-t)^{n-k}dt\\
  &=\binom{n}{k}(\left[ -t^k \cdot \frac{(1-t)^{n-k+1}}{n-k+1} \right]_0^1-\int_0^1 k t^{k-1} \cdot \frac{(1-t)^{n-k+1}}{n-k+1} \,dt)\\
  &=0 + \frac{k}{n-k+1}\int_0^1 \binom{n}{k}t^{k-1}(1-t)^{n-k+1}dt\\
\end{align*}
We repeatedly use the distribution integration formula k times, and it remains:
\[
\int_{0}^{1} \binom{n}{k}t^{k}(1-t)^{n-k}dt = \frac{k!(n-k)!}{(n)!}\int_0^1 t^0 (1-t)^n \, dt
\int_0^1 t^0 (1-t)^n \, dt = \int_0^1 (1-t)^n \, dt. = \frac{1}{n+1}
\]
So, we have:
\[
\int_0^1 b_{n,k}(t) \, dt = \binom{n}{k} \cdot \frac{k! (n-k)!}{(n+1)!} = \frac{1}{n+1}.
\]
The proof is completed.
% ===============================================
\section*{ \center{\normalsize {Acknowledgement}} }
\printbibliography


\end{document}